\задача % 26п.7. Разрезание треугольника
Покажите, как разрезать равносторонний треугольник
\пункт на 9 одинаковых частей;
\пункт на 12 одинаковых частей.
\кзадача

\ВосстановитьГраницы
\vspace{-1mm}
\указание
\пункт Попробуйте вместо разрезания большого треугольника собрать его из 9 маленьких равносторонних треугольников.\\
\пункт Сначала придумайте, как разделить равносторонний треугольник на 3 одинаковые части и как разделить его на 4 одинаковые части.
\куказание

\setcounter{itemnum}{0}

\righttikzw{-85mm}{1mm}{90mm}{
\begin{tikzpicture}[font=\bfseries\fontsize{32}{36}\selectfont]
    \begin{scope}[scale=3.0]
        \node at (0.5,2.4) {\пункт};
        \draw[black, line width=4.2pt]
            (0,0) -- +(0:3) -- +(60:3) -- cycle;
        \draw[red, line width=3.4pt]
            (0,0) ++(0:1)
            -- ++(60:2)
            -- ++(180:1)
            -- ++(-60:2)
            -- ++(60:1)
            -- ++(180:2)
            -- cycle;
    \end{scope}
    \begin{scope}[scale=3.0, xshift=4cm]
        \node at (0.5,2.4) {\пункт};
        \draw[black, line width=4.2pt]
            (0,0) -- +(0:3) -- +(60:3) -- cycle;
        \draw[red, line width=3.4pt]
            (0,0) ++(0:1.5)
            -- ++(60:1.5)
            -- ++(180:1.5)
            -- cycle
            (1.5,0) -- +(150:0.87)
            (0,0) -- ++(30:0.87) -- +(90:0.87)
            (3,0) -- +(150:0.87)
            (1.5,0) -- ++(30:0.87) -- ++(90:0.87) -- +(150:0.87)
            (60:1.5) -- ++(30:0.87) -- +(90:0.87)
            (1.5,0) -- ++(90:0.87) -- +(30:0.87)
            +(0,0) -- +(150:0.87);
    \end{scope}
\end{tikzpicture}}
\ответ
\vspace{29mm}
\кответ
